% Keep styling.sty beside this file. Compile with:
% latexmk -xelatex -synctex=1 template.tex (or, 
% for my purposes, compile to preview with VimtexCompile in neovim)
% if youre using VSC with some built in plugin you'll probably need to add 
% {
%   "latex-workshop.latex.recipe.default": "latexmk (xelatex)",
%   "latex-workshop.latex.autoBuild.run": "onSave",
%   "latex-workshop.view.pdf.viewer": "tab",
%   "latex-workshop.view.pdf.invertMode.enabled": "never"
% }
% also it is necessary to download the right fonts. Just look 
% at error messages to find which ones or something

\documentclass[11pt]{article}
\usepackage{styling} % Dark by default, with the clickable THEME switch.
% \usepackage{graphicx} % Uncomment if adding external images.

\DeclareMathOperator{\Ass}{\operatorname{Ass}}

\title{Assignment 2}
\author{Thobias K. H\o ivik}
\date{For October 27, 2026} % Use \date{} to hide the date.

\begin{document}
\maketitle
\begin{prob*}[1]
    Give in detail the proof of 
    Proposition $8.27$. You must 
    explain every step in detail, 
    including doing Lemma 1.7(b) and 
    Exercise $2.10(a)$. 
    You may take Lemma $2.4$ as given. 
\end{prob*}
\begin{proof}[L{\o}sning]
    Vi ønsker å bevise 
    f{\o}lgende proposisjon, 
    \begin{prop*}[Gathmann Proposition 8.27]
        La $Q_1,\ldots,Q_n$ forme 
        en minimal prim{\ae}r 
        dekomposisjon for 
        et ideal $I$ i en ring $R$, 
        og la $P_i = \sqrt{Q_i}$ for 
        $i = 1,\ldots,n$. 

        Da har vi $\{P_1,\ldots,P_n\} 
        = \Ass(I)$. Spesielt har vi at 
        antall kompnenter i en minimal prim{\ae}r 
        dekomposisjon og deres radikaler 
        ikke er avhengige av den valgte dekomposisjonen. 
    \end{prop*}

    La $R$ v{\ae}re en ring og $I$ et ideal av 
    $R$. 

    Vi antar at $Q_1,\ldots,Q_n$ former en 
    minimal prim{\ae}rdekomposisjon for $I$,
    alts{\aa}
    \[
        I = \bigcap_{i=1}^n Q_i. 
    \]
    Vi setter $P_i = \sqrt{Q_i}$. 

    Siden hver $Q_i$ er prim{\ae}r er hver 
    $P_i$ prim.  
    
    Vi begynner med $P_1,\ldots,P_n \in \Ass(I)$. 

    Fastsett $i$. Vi vil bevise at 
    \[
        P_i \in \Ass(I). 
    \]
    Per definisjon betyr dette at vi m{\aa} finne 
    en $a \in R$ slik at 
    \[
        \sqrt{\{r \in R \bigm| ra \in I\}} 
        = \sqrt{I:a} = P_i. 
    \]

    Minimalitet sier at $Q_i$ ikke er redundant, alts{\aa}
    \[
        \bigcap_{j\neq i} Q_j \not\subseteq Q_i. 
    \]
    Derfor finst det 
    \[
        a \in \bigcap_{j\neq i} Q_j, \ \text{men} \ 
        a \not\in Q_i. 
    \] 

    Siden 
    \[
        I = \bigcap_{j=1}^n Q_j, 
    \]
    trenger vi å bevis at 
    \[
        I : a = \bigcap_{j=1}^n (Q_j : a). 
    \]
    
    Vi observerer at 
    \begin{align*}
        r \in (I : a) 
        &\Leftrightarrow ra \in I \\ 
        &\Leftrightarrow ra \in Q_j \ \forall j \\ 
        &\Leftrightarrow r \in Q_j : a \ \forall j, 
    \end{align*}
    så dette stemmer. 

    N{\aa} ta radikalene, 
    \[
        \sqrt{I:a} = \sqrt{\bigcap_j (Q_j : a)}. 
    \]
    Lemma $1.7(b)$ gir 
    \[
        \sqrt{\bigcap_j (Q_j : a)} = 
        \bigcap_j \sqrt{Q_j : a}. 
    \]

    \begin{proof}[Bevis av Lemma 1.7(b)]
        Vi vil vise at 
        \[
            \sqrt{\bigcap_{j=1}^n  I_j} = 
            \bigcap_{j=1}^n \sqrt{I_j}. 
        \]

        Det holder {\aa} bevise at 
        \[
            \sqrt{I \cap J} = \sqrt I \cap \sqrt J 
        \]
        og konkludere induktivt med noe som
        \begin{align*}
            \sqrt{\bigcap_{j=1}^n I_j} 
            &= \sqrt{
            \underbrace{\left(\bigcap_{j=1}^{n-1} I_j \right)}_{\text{ideal}}
            \cap I_n} \\ 
            &= \sqrt{\bigcap_{j=1}^{n-1} I_j} \cap 
            \sqrt{I_n} \\ 
            &= \sqrt{\bigcap_{j=1}^{n-2}I_j} \cap 
            \sqrt{I_{n-1}} \cap \sqrt{I_n} \\ 
            &\vdots \\ 
            &= \bigcap_{j=1}^n \sqrt{I_j}. 
        \end{align*}

        Viss $a \in \sqrt{I \cap J}$ s{\aa} 
        har vi $a^n \in I \cap J$ for en 
        $n \in \mathbb N$, s{\aa} 
        $a^n \in I$ og $a^n \in J$ og dermed 
        har vi $a \in \sqrt{I} \cap \sqrt{J}$. 

        Anta deretter at $a \in \sqrt{I} \cap \sqrt{J}$. 
        Da har vi $a^n \in I$ og $a^m \in J$ for 
        $n,m \in \mathbb N$ s{\aa} 
        $a^{n+m} = (\underbrace{a^n}_{\in I})(a^m) \in I$
        og $(\underbrace{a^m}_{\in J})(a^n) \in J$ s{\aa} 
        $a^{n+m} \in I \cap J$ og dermed 
        $a \in \sqrt{I \cap J}$. 
    \end{proof}

    N{\aa} gjenst{\aa}r det {\aa} vise at 
    \[
        \bigcap_{j=1}^n \sqrt{Q_j : a} = P_i.
    \]
    Vi valgte $a$ slik at $a \in Q_j$ for $j \neq i$ og 
    $a\neq \in Q_j$ ellers. 

    For $j \neq i$, siden $a \in Q_j$ har vi 
    \[
        Q_j : a = R 
    \]
    siden 
    \[
        Q_j : a = \{r \in R \bigm| ra \in Q_j\}. 
    \]

    For $j = i$, siden $a\not\in Q_j$, 
    har vi at 
    \[
        \sqrt{Q_i : a} = P_i.
    \]

    Vi kan se dette 
    gjennom {\aa} observere at 
    \[
        Q_i \subseteq Q_i : a \subseteq P_i
    \]
    og deretter ta 
    \[
        \sqrt{Q_i} = P_i \subseteq \sqrt{Q_i : a} 
        \subseteq \sqrt{P_i} = P_i. 
    \]

    Da har vi 
    \[
        \sqrt{I:a} = P_i, 
    \]
    og dermed 
    \[
        \{P_1,\ldots,P_n\} \subseteq \Ass(I).
    \]

    N{\aa} mangler det {\aa} bevise at 
    \[
        \Ass(I) \subseteq \{P_1,\ldots,P_n\}. 
    \]

    La $P \in \Ass(I)$ v{\ae}re vilk{\aa}rlig. 
    Per definisjon finnes det $a \in R$ slik at 
    \[
        P = \sqrt{I:a} = \bigcap_{j=1}^n \sqrt{Q_j : a}.
    \]

    Vi beviser n{\aa}  
    Exercise $2.10(a)$ fra Gathmann, 
    \begin{prob*}[Exercise 2.10(a), Gathmann]
        La $J_1,\ldots,J_n$ er idealer og 
        $P$ er prim med 
        \[
            \bigcap_{i=1}^n J_i \subseteq P, 
        \]
        s{\aa} finnes det $i \in [n]$ slik at 
        \[
            J_i \subseteq P. 
        \]
    \end{prob*}
    \begin{proof}[Bevis for $2.10(a)$]
        Anta, for en kontradiksjon, at slik en 
        $i$ ikke finnes, alts{\aa} 
        \[
            J_i \not\subseteq P 
        \]
        for alle $i$. 

        La $a_i \in J_i\setminus P$. 
        Da har vi 
        \[
            \prod_{i=1}^n a_i \in 
            \bigcap{i=1}^n J_i \subseteq P,  
        \]
        og $P$ er prim, s{\aa} 
        for en minst $i$ har vi 
        $a_i \in P$, men vi har antatt 
        $a_i \not\in P$, en kontradiskjon. 
    \end{proof}

    Annvend dette for {\aa} f{\aa} 
    \[
        J_i = \sqrt{Q_i : a} \subseteq P 
    \]
    for minst en $i$. 

    Vi har ogs{\aa} $P \subseteq J_i$ for all 
    $i$ siden $P$ er snittet deres. 

    Dermed har vi 
    \[
        P = J_i = \sqrt{Q_i : a} = P_i 
    \] 
    s{\aa} 
    \[
        \Ass(I) \subseteq \{P_1,\ldots,P_n\}. 
    \]

    $\Ass(I)$ er definert helt med hensyn 
    p{\aa} I, uavhengig av noen dekomposisjon. 

    Definisjonen av en minimal prim{\ae}rdekomposisjon 
    krever at $P_i \neq P_j, \ i \neq j$ som 
    sikrer at ingen ulike prim{\ae}rdekomposisjoner 
    deler samme radikal. 

    Vi f{\aa}r en en-til-en korrespondans mellom 
    $\Ass(I)$ og $P_i$-er. 

\end{proof}

\begin{prob*}[3]
    Skriv den Artinske ringen 
    $k[x,y]/((x-1)^2x,y^2)$ som 
    et produkt av lokale Artinske ringer. 
\end{prob*}
\begin{proof}[Løsning]
    Vi finner de maksimale 
    idealene av $R = k[x,y]/((x-1)^2x, y^2)$. 

    De maksimale idealene av en kvotient ring 
    $R/I$ korresponderer til de maksimale idealene 
    av ringen som inneholder $I$. 

    I dette tilfellet finner vi 
    $\mathfrak m$ maksimale idealer av 
    $k[x,y]$ som inneholder $((x-1)^2x,y^2)$. 

    Viss $(x-1)^2x \in \mathfrak m$ har 
    vi $x \in \mathfrak m$ eller $(x-1)^2 \in \mathfrak m$ 
    som igjen tilsier at $(x-1)\in \mathfrak m$ siden 
    $\mathfrak m$ er prim. 

    Vi har også $y^2 \in \mathfrak m$ så 
    $y \in \mathfrak m$. 

    Da har vi 
    \[
        \mathfrak m_1 = (x,y), \ 
        \mathfrak m_2 = (x-1,y). 
    \]

    Note: Scrap alt som komemr her. 

    Skriv $I = ((x-1)^2x, y^2)$. 

    Vi har 
    \[
        \overline{\mathfrak m_1} 
        \cong \frac{\mathfrak m_1}{I} \ \text{og} \ 
        \overline{\mathfrak m_2} 
        \cong \frac{\mathfrak m_2}{I} 
        \cong k \times k.
    \]
    så vi kan 
    anvende det tredje isomorfi teoremet 
    for å få 
    \[
        R/\overline{\mathfrak m_1} 
        \cong \frac{k[x,y]/I}{\mathfrak m_1/I} 
        \cong \frac{k[x,y]}{(x,y)} \cong k 
    \]
    så $\mathfrak m_1 = (x,y)$ er maksimal. 
    Deretter 
    \[
        R/\overline{\mathfrak m_2} 
        \cong \frac{k[x,y]/I}{\mathfrak m_2/I} 
        \cong \frac{k[x,y]}{(x-1,y)} \cong k 
    \]
    så $\mathfrak m_2 = (x-1,y)$ er maksimal. 

    Vi kan anvende det kinesiske restteoremet for å få 
    \[
        R/\left(\overline{\mathfrak m_1} \cap 
        \overline{\mathfrak m_2}\right) 
        \cong R/(\overline{\mathfrak m_1} \cdot 
        \overline{\mathfrak m_2}) 
        \cong \frac{R}{\overline{\mathfrak m_1}} 
        \cdot \frac{R}{\overline{\mathfrak m_2}} 
    \]
    via den surjektive homomorfien 
    $\varphi : R \to R/\left(\overline{\mathfrak m_1} 
    \cap \overline{\mathfrak m_2}\right)$.  

    Det gjennstår å vise at 
    \[
        \ker \varphi = 0. 
    \]

\end{proof}

\end{document}
